\centerline{\bf \Huge Обратный ход}
\title{Обратный ход}


\problem{0.1} Крестьянин, покупая товары, уплатил первому купцу половину своих денег и ещё 1 рубль; потом уплатил второму купцу половину оставшихся денег да ещё 2 рубля и, наконец, уплатил третьему купцу половину оставшихся да ещё 1 рубль. После этого денег у крестьянина не осталось. Сколько рублей у него было первоначально?

\problem{0.2} У Джона была полная корзина тремпончиков. Сначала он встретил Анну и дал ей половину своих тремпончиков и еще полтремпончика. Потом он встретил Банну и отдал ей половину оставшихся тремпончиков и еще полтремпончика. После того, как он встретил Ванну и снова отдал ей половину тремпончиков и еще полтремпончика, корзина опустела. Сколько тремпончиков было у Джона вначале? (Что такое тремпончики выяснить не удалось, так как к концу задачи их не осталось.)

\problem{0.3} На прямой отметили несколько точек. После этого между каждыми двумя соседними точками добавили по точке. Такую операцию повторили три раза, и в результате на прямой оказалось 65 точек. Сколько точек было вначале?

\problem{1} В 15-этажном доме имеется лифт с двумя кнопками: "+7" и "–9". Можно ли проехать с 3-го этажа на 12-й? 

\problem{2} Дедка вдвое сильнее Бабки, Бабка втрое сильнее Внучки, Внучка вчетверо сильнее Жучки, Жучка впятеро сильнее Кошки, Кошка вшестеро сильнее Мышки. Дедка, Бабка, Внучка, Жучка и Кошка вместе с Мышкой могут вытащить Репку, а без Мышки — не могут. Сколько надо позвать Мышек, чтобы они смогли сами вытащить Репку?

\problem{3} Один Бездельник захотел получить денег и заключил сделку с Чёртом. Теперь каждый раз, когда Бездельник переходит мост через речку, количество имеющихся у него денег удваивается. Но за это он отдаёт Чёрту каждый раз по 24 копейки. Сколько денег было у Бездельника, если он прошёл по мосту 3 раза и деньги у него закончились?

\problem{4} На прямой отметили несколько точек. После этого между каждыми двумя соседними точками отметили ещё по точке. Такое ''уплотнение'' повторили ещё дважды (всего 3 раза). В результате на прямой оказалось отмечено 113 точек. Сколько точек было отмечено первоначально? 

\problem{5} Три мальчика делили 120 фантиков. Сначала Петя дал Ване и Толе столько фантиков, сколько у них было. Затем Ваня дал Толе и Пете столько, сколько у них стало. И наконец, Толя дал Пете и Ване столько, сколько у них к тому моменту имелось. В результате всем досталось поровну. Сколько фантиков было у каждого в начале? 

\newpage

\problem{6} Летит по небу лебедь, а навстречу ему гуси. "Здравствуйте, 100 гусей", - говорит им лебедь, а они ему отвечают:"Нас не 100! А если к нам подлетит ещё столько, сколько нас, и ещё половина, и ещё четверть, и вместе с тобой нас станет 100!". Сколько гусей летело по небу?

\problem{7} Трём братьям дали 24 бублика так, что каждый получил на 3 бублика меньше, чем ему лет. Меньший брат был сообразительным и предложил поменять часть бубликов:"Я, - сказал он, - оставлю половину бубликов, а другую разделю между вами поровну; после этого средний брат также оставит половину бубликов, а другую разделит поровну между мной и старшим братом. В конце старший брат поделит так же." Так они и сделали. Оказалось, что все получили поровну. Сколько лет каждому брату?

\problem{8} По кругу расставлено девять чисел – четыре единицы и пять нулей. Каждую секунду над числами проделывают следующую операцию: между соседними числами ставят ноль, если они различны, и единицу, если они равны; после этого старые числа стирают. Могут ли через некоторое время все числа стать одинаковыми?

 \problem{9} Вовочка подошел к игровому автомату, на экране которого горело число 0. В правилах игры было написано: «На экране показано число очков. Если кинуть монетку в 1 руб., то число очков увеличится на 1. Если кинуть монетку 2 руб., то число очков удвоится. Если набрать 50 очков, то автомат выдаёт приз. А если получилось число, большее 50, то все набранные очки сгорают.» За какое минимальное количество рублей Вовочка сможет получить приз?